\section{Projektbeschreibung, Motivation und Zielgruppe}

\subsection{Ausgangssituation und Projektidee}
Mit einer wachsenden Pokémon-Kartensammlung wird es schwieriger, den Bestand im Blick zu behalten. Sammler möchten wissen, welche Karten sie besitzen, welche Exemplare fehlen und wie oft eine Karte vorhanden ist. Investmentorientierte Nutzer interessieren sich zusätzlich für Marktpreise und Wertentwicklung.

Vergleichbare Anwendungen sind bereits vorhanden. Nach Einschätzung des Projektteams sind viele davon unübersichtlich, wirken veraltet und bieten teure Abomodelle. Daraus entsteht die Motivation, eine Alternative mit einer zeitgemäßen und klar strukturierten Oberfläche zu entwickeln. CardHub verbindet Sammlungsverwaltung und Marktübersicht. Nutzer können Karten suchen, per Kamera erfassen und ihrer Sammlung hinzufügen. Einträge enthalten Anzahl, Kartendaten und verfügbare Preise. Die App ist für Android und iOS geplant und richtet sich an den deutschsprachigen Markt.

\subsection{Motivation und langfristige Vision}
Die Projektidee entstand aus dem eigenen Interesse des Teams an Pokémon-Karten. CardHub soll Erfassung und Verwaltung vereinfachen und Sammlung, Kartendetails und Marktwerte an einem Ort zugänglich machen.

Ein weiteres Ziel ist eine umfassende Übersicht über die Marktentwicklung verschiedener Karten. Nutzer sollen Marktwerte vergleichen und Veränderungen ihres Sammlungswerts nachvollziehen können. Langfristig zeigen Preisverläufe die Entwicklung einzelner Karten. So unterstützt CardHub das Sammeln als Hobby und die Beobachtung von Karten als Investment.

Das Endprodukt umfasst Pokédex, Sammlung, Tausch, Freunde und Profil. Nutzer sollen freigegebene Sammlungen ansehen und Tauschvorschläge austauschen können. Freundesfunktionen, Kartentausch, Gegenangebote, Tauschhistorie, Wunschlisten, Set-Fortschritt und das Erfassen mehrerer Karten in einem Durchgang folgen nach dem MVP.

\subsection{Zielgruppe und Folgen für die Bedienung}
Die primäre Zielgruppe sind deutschsprachige Erwachsene von 18 bis 40 Jahren mit einer physischen Pokémon-Kartensammlung. Die Altersspanne und die folgenden Merkmale sind Planungsannahmen für das Konzept. Hobby-Sammler suchen Überblick und Vollständigkeit. Investmentorientierte Nutzer verfolgen Marktpreise und Wertentwicklung. Deshalb trennt die Oberfläche Bestandsverwaltung und Marktinformationen klar.

Grundlegende Smartphone-Kenntnisse werden vorausgesetzt. Spielregeln und technisches Fachwissen sind nicht erforderlich. Bild, Name, Set und Nummer erleichtern die Zuordnung. Benötigt werden ein Android-Smartphone oder iPhone mit Kamera und für Sammlungsänderungen Internet. Zusätzliche Hardware ist nicht vorgesehen.

Das Konzept geht von sammelinteressierten, ordnungsorientierten und preisbewussten Nutzern aus. Daraus folgen eine übersichtliche Sammlung, eine schnelle Suche und einfache Mengenänderungen.

Unterschiede bei Sehvermögen und Treffsicherheit werden berücksichtigt. Dafür sind anpassbare Schriftgrößen, klare Kontraste und große Bedienelemente vorgesehen. Hinweise werden durch Text und nicht allein durch Farbe vermittelt.

Vorgesehene Situationen sind wiederholte Scans nach dem Öffnen von Boostern, kurze Preischecks unter Zeitdruck und das Ordnen in der Freizeit. Scan und Suche bleiben direkt erreichbar. Nach einem Scan kann die nächste Karte erfasst werden.

Die App soll neben Ausbildung, Studium, Beruf oder Familienaufgaben nutzbar sein. Einkommen und Familienstand werden nicht eingegrenzt. Unterschiedliche Budgets und begrenzte Freizeit sprechen für kurze, unterbrechbare Abläufe und klar bezeichnete Preise mit Datenstand.

